\section{Análisis descriptivo}

El análisis exploratorio de datos permite describir la distribución empírica antes de proponer un modelo probabilístico. En los materiales del curso se enfatiza que cuartiles, IQR, boxplot e histograma permiten estudiar concentración, variabilidad y valores atípicos \cite{curso_eda}. La Tabla \ref{tab:resumen-descriptivo} resume las medidas principales de \texttt{uyu\_por\_brl}; las tablas completas quedan disponibles en \texttt{outputs/tables/}.

\begin{table}[htbp]
\centering
\small
\caption{Resumen descriptivo de la cotización UYU/BRL.}
\label{tab:resumen-descriptivo}
\begin{tabular}{lr@{\hspace{1cm}}lr}
\toprule
Medida & Valor & Medida & Valor \\
\midrule
Media & \MediaCambio{} & Mediana & \MedianaCambio{} \\
Desvío estándar & \DesvioCambio{} & Varianza & \VarianzaCambio{} \\
Q1 & \QunoCambio{} & Q3 & \QtresCambio{} \\
IQR & \IqrCambio{} & Outliers IQR & \OutliersCambio{} \\
Coef. variación & \CoefVariacionCambio{} & Asimetría & \AsimetriaCambio{} \\
\bottomrule
\end{tabular}
\end{table}

La media resume el centro de la distribución, pero no describe por sí sola la variabilidad; por eso se complementa con varianza y desvío estándar, siguiendo el enfoque del curso sobre media y varianza \cite{curso_media_varianza}. El criterio de \(1.5IQR\) no detectó valores atípicos, aunque esto no elimina la presencia de cambios temporales relevantes.

\begin{figure}[htbp]
    \centering
    \includegraphics[width=0.82\textwidth]{serie_uyu_por_brl.png}
    \caption{Evolución temporal diaria de la cotización expresada en pesos uruguayos por real brasileño. Fuente: elaboración propia a partir de Yahoo Finance.}
    \label{fig:serie-cambio}
\end{figure}

La Figura \ref{fig:serie-cambio} evidencia que la cotización no se comporta como una colección de observaciones sin orden. Existen períodos de suba, baja y cambios de nivel. Esta estructura temporal es central: antes de aplicar inferencia clásica, debe considerarse que observaciones próximas en el tiempo pueden estar relacionadas.

\begin{figure}[htbp]
    \centering
    \begin{subfigure}{0.47\textwidth}
        \includegraphics[width=\textwidth]{hist_uyu_por_brl.png}
        \caption{Histograma}
        \label{fig:hist-cambio}
    \end{subfigure}
    \hfill
    \begin{subfigure}{0.47\textwidth}
        \includegraphics[width=\textwidth]{boxplot_uyu_por_brl.png}
        \caption{Boxplot}
        \label{fig:box-cambio}
    \end{subfigure}
    \caption{Distribución empírica y resumen gráfico de la cotización UYU/BRL. Fuente: elaboración propia a partir de Yahoo Finance.}
    \label{fig:dist-box}
\end{figure}

El histograma de la Figura \ref{fig:dist-box} muestra la distribución de frecuencias por clases, tal como se trabaja en el curso para aproximar visualmente la distribución empírica. El boxplot resume mediana, cuartiles y dispersión central. La revisión de retornos y volatilidad se conserva en los archivos generados, pero no se inserta como figura adicional para evitar redundancias en el informe.
